\section{Conclusion}
\label{sec:conclusion}

\emph{Can a single-chip mmWave radar plus \ac{imu} provide usable 6-DoF
state estimation?}  Yes, as a \emph{velocity and orientation}
source with heading aiding: at the live leading edge,
$0.16$\,m/s\,/\,$1.7\degree$ on gentle racing
and $0.28$\,m/s\,/\,$2.6\degree$ on aggressive
racing, with position drifting at $0.3$--$1.3\%$ of distance.
The same configuration is used for all evaluated flights and platforms.
The solver runs faster than its 0.3\,s stride on a laptop CPU,
transfers unchanged to a held-out
flight and to a public dataset, and reports a covariance that is
conservative on racing (\autoref{sec:nees}).  Absolute
position is not observable from Doppler+\ac{imu} and drifts with
distance, and yaw requires a heading sensor.  Through sustained flips
the estimate degrades, reflecting both radar measurement limitations
and the racing-oriented configuration.

Within those structural limits, the answer rested on three things.
First, \emph{characterizing the sensor before designing the
estimator}: the radar's errors have a Gaussian core whose width is a
per-return quantity.  Doppler error from bearing uncertainty grows
linearly with speed, couples through the transverse velocity and
increases off boresight, with a fitted constant of about one
angle-FFT bin.  The errors also
include a component shared by all returns of a frame, an outlier
remainder, and aliased returns with elevated bearing-related residuals
even after Doppler unwrapping.
Second, \emph{deriving
the radar noise weights and the radar and \ac{imu} stream weights
from those measurements}, and
the bias priors and the staleness of the marginalization prior from
measurements of their own, so that the reported covariance is a prediction tested against ground truth (\autoref{sec:nees}).  Third, \emph{a fixed-lag smoother} that
carries past information through an approximate marginalization prior, built from exactly the factors that leave the window,
applied
without a global scale on every regime.  The largest gains
came from measuring and re-weighting the existing measurements.

\emph{Future directions} follow from the limitations.  More precise
per-return radar timing could help model intra-frame motion during
backflips.  The noise floor $\sigma_0$ and the frame-shared term $\sigma_c$
could be adapted in flight, but estimating the shared component
requires separating measurement error from trajectory error.  Repeated recordings per
regime would quantify between-flight variability.  An absolute-height anchor
such as a floor-plane factor~\cite{yang2025gorio} could be gated on the
frames that observe the intermittently visible floor.
